\subsubsection{Linear Differential Operator}

\paragraph{TheundefinedFormalundefinedLinearundefinedDifferentialundefinedOperator}

Weundefineddefineundefinedtheundefineddifferentianundefinedopeartorundefinedas:

\begin{mystthmdefinition}{DifferentiationundefinedOpeartor}{definition-lineardifferentialoperator-1}Letundefined$D:L^2[a,b]\to L^2[a,b]$undefinedbeundefinedtheundefinedmappingundefineddefinedundefinedby:
\begin{equation}
D(u)=\frac{d}{dx}u
\end{equation}
forundefinedallundefined$u\in L^2[a,b]$

\end{mystthmdefinition}
Nowundefinedweundefinedcanundefineddefineundefinedtheundefined\textit{formalundefinedlinearundefinedoperator}undefinedbyundefinedsomeundefinedopeartorundefinedalgebra:

\begin{mystthmdefinition}{Formalundefinedlinearundefineddifferentialundefinedoperator}{definition-lineardifferentialoperator-2}Letundefined$p(x)$undefinedbeundefinedanyundefinedpolynomialundefinedoverundefined$\mathbb{C}$undefinedofundefinedpositiveundefineddegree:
\begin{equation}
p(x)=a_\nu x^\nu
\end{equation}
whereundefined$a_{\nu}$undefinedareundefinedfunctionsundefinedofundefined$x$.undefinedDefineundefined$L:=p(D)$undefinedasundefinedtheundefined\textbf{formalundefinedlinearundefineddifferentialundefinedoperator}undefinedthatundefinedcorrespondundefinedtoundefinedtheundefined\textbf{auxillaryundefinedpolynomialundefined$p$}:
\begin{equation}
L=a_\nu D^{\nu}
\end{equation}
Theundefined\textbf{order}undefinedofundefined$L$undefinedisundefinedgivenundefinedbyundefinedtheundefinedorderundefinedofundefined$p$.

\end{mystthmdefinition}
Weundefinedcallundefineditundefined\textit{formal}undefinedsinceundefinedweundefinedwantundefinedtoundefinedcontrastundefinedwithundefined\textit{concrete}undefinedlinearundefineddifferentialundefinedopeartorundefinedwhichundefinedisundefinedaundefinedformalundefinedopeartorundefinedendorsedundefinedwithundefinedboundaryundefinedconditions(Thisundefinedconventionundefinedisundefinedusedundefinedinundefined\textit{MathematicalundefinedPhysicsundefinedbyundefinedMichaelundefinedStoneundefinedandundefinedPaulundefinedGoldbart}).undefinedWeundefinedcanundefinedthinkundefinedofundefinedtheundefineddifferenceundefinedbetweenundefinedthemundefinedas:

\begin{itemize}
\item Theundefined\textbf{nullspaceundefinedofundefinedformalundefinedopeartor}undefinedisundefinedtheundefinedspaceundefinedofundefinedgeneralundefinedsolutionundefinedofundefinedhomogenousundefineddiffferentialundefinedequation
\item Theundefined\textbf{nullspaceundefinedofundefinedconcreteundefinedoperator}undefinedisundefinedtheundefinedspaceundefinedofundefinedparticularundefinedsolutionundefinedtoundefinedhomogenousundefineddifferentialundefinedequationundefinedwithundefinedboundaryundefinedconditions
\end{itemize}

We'llundefinedexploreundefinedtheundefinednullspaceundefinedofundefinedformalundefinedoperatorsundefinedwithundefinedaundefinedbitundefinedmoreundefineddetailundefinednow.

\subparagraph{Nullspaceundefinedofundefinedformalundefinedoperator}

Weundefinedfirstundefinedobserveundefinedsomethingundefinedveryundefinedniceundefinedaboutundefinedsolutionsundefinedtoundefinedhomogenousundefineddifferentialundefinedequations:undefinedtheyundefinedbelongundefinedtoundefined$C^{\infty}$(infinitelyundefineddifferentiable).

\subparagraph{TheundefinedBootstrappingundefinedLemma}

\begin{mystthmlemma}{}{lemma-lineardifferentialoperator-3}Letundefined$L$undefinedbeundefinedaundefinedformalundefinedoperatorundefinedofundefinedorderundefined$n$.undefinedSupposeundefinedtheundefinedcoefficientsundefinedofundefinedtheundefinedauxillaryundefinedpolynomialundefinedareundefined$k$undefinedtimesundefinedcontinouslyundefineddifferentiable($a_{\nu}\in C^k$)undefinedthenundefined$\forall y\in \mathrm{Null}(L)$,undefined$y$undefinedisundefinedatundefinedleastundefined$k+n$undefinedtimesundefinedcontiouslyundefineddifferentiable.

\end{mystthmlemma}
\begin{mystthmproof}{}{proof-lineardifferentialoperator-4}Supposeundefined$y\in \mathrm{Null}(L)$,undefinednoteundefinedtriviallyundefined$y\in C^n$.undefinedThen:
\begin{equation}
\begin{aligned}
a_0y+a_1y^{(1)}+a_2y^{(2)}...+a_n y^{(n)} &=0 \\
y^{(n)} &= -(b_0y+b_1y^{(1)}+b_2y^{(2)}...\underbrace{b_{n-1}{y}^{n-1}}_{\in C^{\min(k, 1)}})
\end{aligned}
\end{equation}
Whereundefined$b_{\nu}=a_{\nu}/a_{n}$.undefinedNowundefined$b_{\nu}\in C^k$undefinedbyundefinedclosednessundefinedunderundefinedpolynomialundefineddivision.undefinedAlsoundefined$y^{(\nu)} \in C^{n -\nu}$undefinedhenceundefinedtheundefinedR.H.Sundefinedisundefinedatundefinedleastundefinedoneundefinedtimeundefinedcontinouslyundefineddifferentiable(theundefinedbottleundefinedneckundefinedbeingundefined$a_{n -1}y^{(n -1)} \in C^{\min(k, 1)}$).undefinedThereforeundefined$y^{(n)}\in C^1$.undefinedButundefinedthisundefinedimpliesundefined$y\in C^{n+1}$undefinedhenceundefinedweundefinedrepeatundefinedtheundefinedaboveundefinedargumentundefinedtoundefinedgetundefined$y^{(\nu)}\in C^{n+1 -\nu}$.undefinedTheundefinedR.H.Sundefinedhasundefinedtheundefinedbottleneckundefined$a_{n -1}y^{(n -1)}\in C^{\min(k, n+1 -(n -1))}=C^{\min(k, 2)}$.undefinedWeundefinedcanundefinedrepeatundefinedthisundefinedprocedureundefined$m$undefinedtimesundefinedandundefinedyieldundefinedthatundefined$y^{(n)}\in C^{\min(k, m)}$.undefinedThisundefinedmeansundefinedtheundefinedbestundefinedweundefinedcanundefinedgetundefinedisundefined$y^{(n)}\in C^k$undefinedorundefined$y\in C^{n+k}$(afterundefinedwhichundefinedtheundefinedcoeficientundefinedbecomesundefinedtheundefinedbottleneck).

\end{mystthmproof}
\begin{mystthmcorollary}{}{corollary-lineardifferentialoperator-5}Ifundefinedtheundefinedcoefficientsundefinedareundefinedinundefined$C^{\infty}$undefinedthenundefined$y\in \mathrm{Null}(L)$undefinedareundefinedinundefined$C^{\infty}$.

\end{mystthmcorollary}
\subparagraph{Nullspaceundefinedofundefined$L$undefinedisundefinedanundefined$n$-dimensionalundefinedsubspaceundefinedofundefined$C^{\infty}$}

\paragraph{LinearundefinedIndependence:undefinedWronskian}

\paragraph{Solvingundefined$Ly=0$undefinedwithundefinedconstantundefinedcoefficients}

\paragraph{$Ly=0$undefinedwithundefinedgeneralundefined$a_\nu(x)$}

\begin{mystexercise}{1-dimensionalundefinedScatteringundefinedProblem}{1dscat}\label{1dscat}Considerundefinedtheundefined1dundefinedSchrodingerundefinedequation:
\begin{equation}
-\frac{d^2\psi}{dx^2}+V\psi = E\psi
\end{equation}
Thisundefinedisundefinedanundefinedeigenvalueundefinedproblemundefinedofundefinedtheundefinedoperatorundefined$H= -\frac{d^2}{dx^2}+V$:
\begin{equation}
H\psi = E\psi
\end{equation}
Whereundefined$V$undefinedhasundefinedfiniteundefinedsupportundefined$[ -a, a], a\in \mathbb{R}$.undefinedDenote:

\begin{enumerate}
\item $L:=\{x\in \mathbb{R}:x< -a\}$
\item $R:=\{x\in \mathbb{R}:x>a\}$
\end{enumerate}

\textbf{Problems}:

\begin{enumerate}
\item Showundefinedthatundefinedinundefined$L,R$,undefinedtheundefinedgeneralundefinedsolutionundefinedtoundefinedtheundefinedformalundefinedoperatorundefined$H -E$undefinedwhereundefined$E=k^2$undefinedis:

\begin{equation}
\begin{aligned}
\psi_k &=A\exp(ikx)+B\exp(-ikx) \quad x\in L  \\
\psi_k &=C\exp(ikx)+D\exp(-ikx) \quad x\in R
\end{aligned}
\end{equation}


\item Nowundefinedconsiderundefinedtwoundefinedasymptoticundefinedboundaryundefinedcondition:

\begin{equation}
\begin{cases}
\Pi_L: D=0, x\to \infty \\
\Pi_R: C=0, x\to -\infty
\end{cases}
\end{equation}


whereundefined$\Pi_L$undefineddepictsundefinedaundefinedleft-incidentundefinedwaveundefinedwithundefinednoundefinedincomingundefinedwaveundefinedfromundefinedtheundefinedrightundefinedandundefined$\Pi_R$undefinedisundefinedtheundefinedopposite.undefinedShowundefinedthatundefinedtheundefinedsolution(inundefined$L,R$)undefinedforundefinedtheundefinedBCundefined$\Pi_1$undefinedandundefined$k>0$undefinedis:

\begin{equation}
\begin{aligned}
\psi_k^L(x) = \begin{cases}
\exp(ikx)+ r_L(k)\exp(-ikx) & x\in L \\
t_L(k)\exp(ikx) & x\in R
\end{cases}
\end{aligned}
\end{equation}


Similarily,undefinedforundefined$\Pi_2$,undefined$k<0$:

\begin{equation}
\begin{aligned}
\psi_k^R(x)=\begin{cases}
t_R(k)\exp(ikx)& x\in L \\
\exp(ikx) + r_R(k)\exp(-ikx) & x\in R
\end{cases}
\end{aligned}
\end{equation}


whereundefined$t, r$undefineddenotesundefinedtheundefined\textit{transmission}undefinedandundefined\textit{reflection}undefinedcoefficient.


\item Noteundefinedthatundefined$\psi_k^*$undefinedisundefinedalsoundefinedaundefinedsolution.undefinedUseundefinedpropertiesundefinedofundefinedWronskianundefinedtoundefinedshwoundefinedthat:

\begin{equation}
|t_{L/R}|^2 + |r_{L/R}|^2=1
\end{equation}
\end{enumerate}

\end{mystexercise}
\begin{mystsolution}{SolutionundefinedtoundefinedExercise~\ref{1dscat}}\begin{enumerate}
\item Fixundefinedaundefined$k$.undefinedConsiderundefinedtheundefinedbasisundefined$\psi_k^L, {\psi_k^L}^*$.undefinedFormundefinedtheundefinedWronskianundefinedinundefined$L$:

\begin{equation}
\begin{aligned}
W_L=\begin{vmatrix}
\exp(ikx)+r_L\exp(-ikx) &   \exp(-ikx)+r_L^*\exp(ikx) \\
ik\exp(ikx)-ikr_L\exp(-ikx) & -ik\exp(-ikx)+ikr_L^*\exp(ikx)
\end{vmatrix}
\end{aligned}
\end{equation}


Evaluateundefinedtheundefineddeterminant:

\begin{equation}
\begin{aligned}
W_L &= (-ik)\cancel{-ikr_L\exp(-2ikx) + ikr_L\exp(2ikx)}+ik|r_L|^2  \\
&-[ik\cancel{-ikr_K\exp(-2ikx)+ikr_L\exp(2ikx)}-ik|r_L|^2]  \\
&= 2ik(|r_L|^2-1)
\end{aligned}
\end{equation}


NextundefinedformundefinedWronskianundefinedinundefined$R$:

\begin{equation}
\begin{aligned}
W_R = \begin{vmatrix}
t_L\exp(ikx) & t_L\exp(-ikx) \\
ikt_L\exp(ikx) & -ikt_L\exp(-ikx) 
\end{vmatrix} = -2ik|t_L|^2
\end{aligned}
\end{equation}


Sinceundefinedweundefinedwouldundefinedexpectundefined$W_R=W_L=0$,undefinedhenece:

\begin{equation}
\begin{aligned}
W_R-W_L&=0 \\
2ik(|r_L|^2+|t_L|^2-1) &=0
\end{aligned}
\end{equation}


Sinceundefined$k>0$undefinedweundefinedgetundefinedasundefineddesired.
\end{enumerate}

4.Fixundefined$k$.undefinedWeundefinedrepeatundefinedtheundefinedsameundefinedprocedureundefinedwithundefined$\{\psi_k^L,\psi_{( -k)}^R\}$.undefinedWeundefinedevaluateundefined$W_L$:
\begin{equation}
\begin{aligned}
W_L &= \begin{vmatrix}
\exp(ikx)+r_L(k)\exp(-ikx) & t_R(-k)\exp(-ikx) \\
ik\exp(ikx)-ikr_L(k)\exp(-ikx) & -ikt_R(-k)\exp(-ikx)
\end{vmatrix} \\
&= -2ikt_R(-k)
\end{aligned}
\end{equation}
Similarily,
\begin{equation}
W_R = -2ikt_L(k)
\end{equation}
Again,undefined$W_L=W_R=0$undefinedsoundefinedweundefinedgetundefinedasundefineddesired.

\end{mystsolution}